% ============================================================
\chapter{$\sigma$-Alg\`ebres et Espaces Mesurables}
\label{ch:sigma_algebres}
% ============================================================

\`A la fin du XIX\ieme{} si\`ecle, les math\'ematiciens se heurtaient \`a un
probl\`eme fondamental~: comment d\'efinir rigoureusement la \og taille\fg{}
d'un ensemble~? L'int\'egrale de Riemann, qui avait si bien servi pendant
des d\'ecennies, commen\c{c}ait \`a montrer ses limites face aux fonctions
pathologiques d\'ecouvertes par Dirichlet, Weierstrass et d'autres. Henri
Lebesgue, dans sa th\`ese de 1902, proposa une r\'evolution~: au lieu de
d\'ecouper l'axe des $x$ en petits intervalles, il d\'ecoupait l'axe des
$y$. Mais pour que cette id\'ee fonctionne, il fallait d'abord r\'epondre
\`a une question pr\'ealable~: \`a quels ensembles peut-on attribuer une
mesure de mani\`ere coh\'erente~?

La r\'eponse --- et c'est l'objet de ce premier chapitre --- est qu'il
faut se restreindre \`a une collection d'ensembles poss\'edant de bonnes
propri\'et\'es de cl\^oture~: une $\sigma$-alg\`ebre, ou \emph{tribu}. Ce
concept, qui semble au premier abord purement technique, est en r\'ealit\'e
la cl\'e de vo\^ute de toute la th\'eorie de la mesure, de l'int\'egration, et
par extension des probabilit\'es modernes.

\section{Introduction}

La premi\`ere \'etape dans la construction de la th\'eorie de la mesure consiste \`a
identifier les ensembles auxquels on souhaite attribuer une \og taille \fg. On ne
peut pas, en g\'en\'eral, mesurer \emph{tous} les sous-ensembles d'un ensemble donn\'e
de mani\`ere coh\'erente (nous verrons au chapitre~\ref{ch:lebesgue} l'existence
d'ensembles non mesurables au sens de Lebesgue). Il faut donc se restreindre \`a une
collection d'ensembles suffisamment riche pour \^etre utile, mais suffisamment
structur\'ee pour que la mesure soit bien d\'efinie. Cette collection est une
$\sigma$-alg\`ebre (ou \emph{tribu}).

\section{D\'efinitions fondamentales}

\begin{definition}[$\sigma$-alg\`ebre]
Soit $X$ un ensemble non vide. Une \emph{$\sigma$-alg\`ebre} (ou \emph{tribu}) sur $X$
est une collection $\mathcal{A} \subset \mathcal{P}(X)$ de sous-ensembles de $X$
v\'erifiant~:
\begin{enumerate}
    \item[(i)] $X \in \mathcal{A}$\,;
    \item[(ii)] si $A \in \mathcal{A}$, alors $A^c = X \setminus A \in \mathcal{A}$\,;
    \item[(iii)] si $(A_n)_{n \in \N}$ est une suite d'\'el\'ements de $\mathcal{A}$,
    alors $\bigcup_{n=0}^{\infty} A_n \in \mathcal{A}$.
\end{enumerate}
Le couple $(X, \mathcal{A})$ est appel\'e \emph{espace mesurable}.
\end{definition}

\begin{remarque}
La condition (iii) porte sur les r\'eunions \emph{d\'enombrables}. C'est la
diff\'erence essentielle avec la notion d'\emph{alg\`ebre} (ou alg\`ebre de Boole),
qui ne requiert la stabilit\'e que par r\'eunions \emph{finies}.
\end{remarque}

\begin{proposition}[Propri\'et\'es \'el\'ementaires]
\label{prop:sigma_props}
Soit $\mathcal{A}$ une $\sigma$-alg\`ebre sur $X$. Alors~:
\begin{enumerate}
    \item $\varnothing \in \mathcal{A}$\,;
    \item $\mathcal{A}$ est stable par intersection d\'enombrable~: si
    $(A_n)_{n \in \N} \subset \mathcal{A}$, alors
    $\bigcap_{n=0}^{\infty} A_n \in \mathcal{A}$\,;
    \item $\mathcal{A}$ est stable par diff\'erence~: si $A, B \in \mathcal{A}$,
    alors $A \setminus B \in \mathcal{A}$\,;
    \item $\mathcal{A}$ est stable par r\'eunion et intersection finies.
\end{enumerate}
\end{proposition}

\begin{proof}
(1) $\varnothing = X^c \in \mathcal{A}$ par (i) et (ii).

(2) Par les lois de De Morgan~:
$\bigcap_{n} A_n = \left(\bigcup_{n} A_n^c\right)^c$. Comme chaque
$A_n^c \in \mathcal{A}$ par (ii), la r\'eunion $\bigcup_n A_n^c \in \mathcal{A}$
par (iii), et son compl\'ementaire est dans $\mathcal{A}$ par (ii).

(3) $A \setminus B = A \cap B^c$, qui est dans $\mathcal{A}$ par (ii) et (2).

(4) Cas particulier de (iii) et (2) avec $A_n = \varnothing$ pour $n$ assez grand.
\end{proof}

\begin{definition}[Alg\`ebre]
Une collection $\mathcal{A}_0 \subset \mathcal{P}(X)$ est une \emph{alg\`ebre} sur
$X$ si elle v\'erifie (i), (ii) et la stabilit\'e par r\'eunion \emph{finie}~:
si $A, B \in \mathcal{A}_0$, alors $A \cup B \in \mathcal{A}_0$.
\end{definition}

\begin{exemple}
Voici quelques exemples fondamentaux de $\sigma$-alg\`ebres~:
\begin{enumerate}
    \item La \emph{tribu triviale} $\{\varnothing, X\}$ est la plus petite
    $\sigma$-alg\`ebre sur $X$.
    \item La \emph{tribu discr\`ete} $\mathcal{P}(X)$ est la plus grande
    $\sigma$-alg\`ebre sur $X$.
    \item Pour tout $A \subset X$, $\mathcal{A} = \{\varnothing, A, A^c, X\}$ est
    une $\sigma$-alg\`ebre.
    \item Si $X$ est non d\'enombrable, la collection des ensembles d\'enombrables
    ou \`a compl\'ementaire d\'enombrable forme une $\sigma$-alg\`ebre, appel\'ee
    \emph{tribu codénombrable}.
\end{enumerate}
\end{exemple}

\section{$\sigma$-Alg\`ebre engendr\'ee}

En pratique, on ne d\'efinit pas une $\sigma$-alg\`ebre en listant tous ses
\'el\'ements, mais en la \og engendrant \fg \`a partir d'une collection plus petite.

\begin{theoreme}[$\sigma$-alg\`ebre engendr\'ee]
\label{thm:sigma_engendree}
Soit $\mathcal{E} \subset \mathcal{P}(X)$ une collection quelconque de
sous-ensembles de $X$. Il existe une plus petite $\sigma$-alg\`ebre contenant
$\mathcal{E}$, appel\'ee la \emph{$\sigma$-alg\`ebre engendr\'ee par $\mathcal{E}$}
et not\'ee $\sigma(\mathcal{E})$. Elle est donn\'ee par~:
\[
\sigma(\mathcal{E}) = \bigcap_{\substack{\mathcal{A} \text{ $\sigma$-alg\`ebre} \\
\mathcal{E} \subset \mathcal{A}}} \mathcal{A}.
\]
\end{theoreme}

\begin{proof}
L'intersection est bien d\'efinie car au moins $\mathcal{P}(X)$ est une
$\sigma$-alg\`ebre contenant $\mathcal{E}$, donc la famille n'est pas vide.

V\'erifions que $\sigma(\mathcal{E})$ est une $\sigma$-alg\`ebre~:
\begin{itemize}
    \item $X$ appartient \`a toute $\sigma$-alg\`ebre, donc $X \in \sigma(\mathcal{E})$.
    \item Si $A \in \sigma(\mathcal{E})$, alors $A$ est dans toute $\sigma$-alg\`ebre
    $\mathcal{A} \supset \mathcal{E}$, donc $A^c$ l'est aussi, donc
    $A^c \in \sigma(\mathcal{E})$.
    \item De m\^eme pour les r\'eunions d\'enombrables.
\end{itemize}
Enfin, $\sigma(\mathcal{E})$ est la \emph{plus petite} car elle est contenue dans
toute $\sigma$-alg\`ebre contenant $\mathcal{E}$ (par d\'efinition comme intersection).
\end{proof}

\begin{remarque}
La construction par intersection est non constructive~: elle ne permet pas de
d\'ecrire explicitement les \'el\'ements de $\sigma(\mathcal{E})$. En g\'en\'eral,
$\sigma(\mathcal{E})$ est beaucoup plus grande que $\mathcal{E}$ et sa structure
est complexe.
\end{remarque}

\section{La tribu bor\'elienne}

\begin{definition}[Tribu bor\'elienne]
Soit $(X, \tau)$ un espace topologique. La \emph{tribu bor\'elienne} de $X$, not\'ee
$\mathcal{B}(X)$, est la $\sigma$-alg\`ebre engendr\'ee par les ouverts de $X$~:
\[
\mathcal{B}(X) = \sigma(\tau).
\]
Les \'el\'ements de $\mathcal{B}(X)$ sont appel\'es les \emph{ensembles bor\'eliens}
(ou \emph{bor\'eliens}) de $X$.
\end{definition}

\begin{theoreme}[G\'en\'erateurs de la tribu bor\'elienne de $\R$]
\label{thm:generateurs_boreliens}
La tribu bor\'elienne $\mathcal{B}(\R)$ est engendr\'ee par chacune des collections
suivantes~:
\begin{enumerate}
    \item les ouverts de $\R$\,;
    \item les ferm\'es de $\R$\,;
    \item les intervalles ouverts $(a, b)$, $a < b$\,;
    \item les intervalles semi-ouverts $[a, b)$, $a < b$\,;
    \item les intervalles semi-ouverts $(a, b]$, $a < b$\,;
    \item les demi-droites $(-\infty, a)$, $a \in \R$\,;
    \item les demi-droites $(-\infty, a]$, $a \in \R$\,;
    \item les demi-droites $(a, +\infty)$, $a \in \R$.
\end{enumerate}
\end{theoreme}

\begin{proof}
Notons $\mathcal{B}$ la tribu bor\'elienne et $\sigma_i$ la $\sigma$-alg\`ebre
engendr\'ee par la $i$-\`eme collection. Il suffit de montrer les inclusions
circulaires.

$\sigma_6 \subset \sigma_7$~: $(-\infty, a) = \bigcup_{n=1}^{\infty}
(-\infty, a - 1/n]$ donc $(-\infty, a) \in \sigma_7$.

$\sigma_7 \subset \sigma_6$~: $(-\infty, a] = \bigcap_{n=1}^{\infty}
(-\infty, a + 1/n)$ donc $(-\infty, a] \in \sigma_6$.

$\sigma_3 \subset \sigma_6$~: $(a, b) = (-\infty, b) \setminus (-\infty, a]
= (-\infty, b) \cap (a, +\infty)$. Or $(a, +\infty) = (-\infty, a]^c \in \sigma_6$
(car $\sigma_6 = \sigma_7$ comme montr\'e ci-dessus).

$\sigma_1 \subset \sigma_3$~: tout ouvert de $\R$ est r\'eunion d\'enombrable
d'intervalles ouverts (propri\'et\'e de Lindel\"of).

Comme $\sigma_1 = \mathcal{B}$ par d\'efinition et que
$\sigma_1 \subset \sigma_3 \subset \sigma_6 \subset \mathcal{B}$, on a \'egalit\'e.
Les autres cas se traitent de mani\`ere analogue.
\end{proof}

\begin{formules}[title=\textbf{G\'en\'erateurs courants de $\mathcal{B}(\R)$}]
\[
\mathcal{B}(\R) = \sigma\bigl(\{(a,b) : a < b\}\bigr)
= \sigma\bigl(\{(-\infty, a) : a \in \R\}\bigr)
= \sigma\bigl(\{(-\infty, a] : a \in \R\}\bigr).
\]
\end{formules}

\section{Tribu bor\'elienne de $\R^n$}

\begin{definition}
La tribu bor\'elienne de $\R^n$ est $\mathcal{B}(\R^n) = \sigma(\tau_n)$, o\`u
$\tau_n$ est la topologie usuelle de $\R^n$.
\end{definition}

\begin{proposition}
On a $\mathcal{B}(\R^n) = \sigma(\mathcal{R}_n)$, o\`u $\mathcal{R}_n$ est la
collection des pav\'es ouverts
$\prod_{i=1}^n (a_i, b_i)$ avec $a_i < b_i$ pour tout $i$.
\end{proposition}

\begin{proof}
Tout ouvert de $\R^n$ est r\'eunion d\'enombrable de pav\'es ouverts \`a sommets
rationnels (car $\Q^{2n}$ est d\'enombrable), donc
$\sigma(\tau_n) \subset \sigma(\mathcal{R}_n)$. R\'eciproquement, chaque pav\'e
ouvert est un ouvert de $\R^n$, donc $\sigma(\mathcal{R}_n) \subset \sigma(\tau_n)$.
\end{proof}

\section{$\sigma$-Alg\`ebre produit}

\begin{definition}[$\sigma$-alg\`ebre produit]
Soient $(X_1, \mathcal{A}_1)$ et $(X_2, \mathcal{A}_2)$ deux espaces mesurables.
La \emph{$\sigma$-alg\`ebre produit} est~:
\[
\mathcal{A}_1 \otimes \mathcal{A}_2 = \sigma\bigl(\{A_1 \times A_2 :
A_1 \in \mathcal{A}_1,\, A_2 \in \mathcal{A}_2\}\bigr).
\]
Plus g\'en\'eralement, pour une famille $(X_i, \mathcal{A}_i)_{i \in I}$~:
\[
\bigotimes_{i \in I} \mathcal{A}_i = \sigma\bigl(\{\pi_i^{-1}(A_i) :
i \in I,\, A_i \in \mathcal{A}_i\}\bigr),
\]
o\`u $\pi_i : \prod_{j \in I} X_j \to X_i$ est la projection canonique.
\end{definition}

\begin{theoreme}
\label{thm:borel_produit}
Si $X_1$ et $X_2$ sont des espaces m\'etriques s\'eparables, alors
\[
\mathcal{B}(X_1) \otimes \mathcal{B}(X_2) = \mathcal{B}(X_1 \times X_2).
\]
En particulier, $\mathcal{B}(\R^n) = \mathcal{B}(\R) \otimes \cdots \otimes
\mathcal{B}(\R)$ ($n$ fois).
\end{theoreme}

\begin{proof}
Soit $U$ un ouvert de $X_1 \times X_2$. Pour tout $(x_1, x_2) \in U$, il existe
des ouverts $V_1 \ni x_1$ et $V_2 \ni x_2$ tels que $V_1 \times V_2 \subset U$.
Par s\'eparabilit\'e, on peut choisir $V_1$ et $V_2$ dans des bases d\'enombrables
$\mathcal{B}_1$ et $\mathcal{B}_2$ de $X_1$ et $X_2$ respectivement. Donc
$U = \bigcup_{k} V_1^{(k)} \times V_2^{(k)}$, r\'eunion d\'enombrable de rectangles
mesurables. Ceci montre $\mathcal{B}(X_1 \times X_2) \subset
\mathcal{B}(X_1) \otimes \mathcal{B}(X_2)$.

R\'eciproquement, si $A_1 \in \mathcal{B}(X_1)$, alors $\pi_1^{-1}(A_1) =
A_1 \times X_2$ est bor\'elien dans $X_1 \times X_2$ car $\pi_1$ est continue.
De m\^eme pour $\pi_2$. Les ensembles $A_1 \times A_2 = \pi_1^{-1}(A_1) \cap
\pi_2^{-1}(A_2)$ sont bor\'eliens, d'o\`u l'inclusion r\'eciproque.
\end{proof}

\section{Sous-$\sigma$-alg\`ebre et trace}

\begin{definition}[Trace]
Soit $(X, \mathcal{A})$ un espace mesurable et $Y \subset X$. La \emph{trace} de
$\mathcal{A}$ sur $Y$ est~:
\[
\mathcal{A}|_Y = \{A \cap Y : A \in \mathcal{A}\}.
\]
C'est une $\sigma$-alg\`ebre sur $Y$.
\end{definition}

\begin{proposition}
Si $\mathcal{A} = \sigma(\mathcal{E})$, alors
$\mathcal{A}|_Y = \sigma(\mathcal{E}|_Y)$, o\`u
$\mathcal{E}|_Y = \{E \cap Y : E \in \mathcal{E}\}$.
\end{proposition}

\begin{proof}
Posons $\mathcal{F} = \{A \subset X : A \cap Y \in \sigma(\mathcal{E}|_Y)\}$.
On v\'erifie que $\mathcal{F}$ est une $\sigma$-alg\`ebre contenant $\mathcal{E}$,
donc $\sigma(\mathcal{E}) \subset \mathcal{F}$, ce qui donne
$\mathcal{A}|_Y \subset \sigma(\mathcal{E}|_Y)$. L'inclusion r\'eciproque est
imm\'ediate.
\end{proof}

\section{Applications mesurables}

\begin{definition}[Application mesurable]
Soient $(X, \mathcal{A})$ et $(Y, \mathcal{B})$ deux espaces mesurables. Une
application $f : X \to Y$ est dite \emph{$(\mathcal{A}, \mathcal{B})$-mesurable}
(ou simplement \emph{mesurable}) si~:
\[
\forall\, B \in \mathcal{B}, \quad f^{-1}(B) \in \mathcal{A}.
\]
\end{definition}

\begin{proposition}[Crit\`ere de mesurabilit\'e par g\'en\'erateurs]
\label{prop:critere_mesurabilite}
Soit $\mathcal{B} = \sigma(\mathcal{E})$. Alors $f : (X, \mathcal{A}) \to
(Y, \mathcal{B})$ est mesurable si et seulement si~:
\[
\forall\, E \in \mathcal{E}, \quad f^{-1}(E) \in \mathcal{A}.
\]
\end{proposition}

\begin{proof}
Le sens direct est \'evident. Pour la r\'eciproque, posons
$\mathcal{C} = \{B \subset Y : f^{-1}(B) \in \mathcal{A}\}$.
V\'erifions que $\mathcal{C}$ est une $\sigma$-alg\`ebre~:
\begin{itemize}
    \item $f^{-1}(Y) = X \in \mathcal{A}$, donc $Y \in \mathcal{C}$.
    \item Si $B \in \mathcal{C}$, alors $f^{-1}(B^c) = f^{-1}(B)^c \in \mathcal{A}$,
    donc $B^c \in \mathcal{C}$.
    \item Si $(B_n) \subset \mathcal{C}$, alors
    $f^{-1}(\bigcup_n B_n) = \bigcup_n f^{-1}(B_n) \in \mathcal{A}$, donc
    $\bigcup_n B_n \in \mathcal{C}$.
\end{itemize}
Par hypoth\`ese, $\mathcal{E} \subset \mathcal{C}$, donc
$\mathcal{B} = \sigma(\mathcal{E}) \subset \mathcal{C}$, ce qui signifie que $f$
est mesurable.
\end{proof}

\begin{corollaire}
Une application $f : \R \to \R$ est bor\'elienne (i.e.\
$(\mathcal{B}(\R), \mathcal{B}(\R))$-mesurable) si et seulement si
$f^{-1}((-\infty, a)) \in \mathcal{B}(\R)$ pour tout $a \in \R$.
\end{corollaire}

\begin{proposition}[Composition]
Si $f : (X, \mathcal{A}) \to (Y, \mathcal{B})$ et
$g : (Y, \mathcal{B}) \to (Z, \mathcal{C})$ sont mesurables, alors
$g \circ f : (X, \mathcal{A}) \to (Z, \mathcal{C})$ est mesurable.
\end{proposition}

\begin{proof}
Pour tout $C \in \mathcal{C}$, $(g \circ f)^{-1}(C) = f^{-1}(g^{-1}(C))$.
Comme $g$ est mesurable, $g^{-1}(C) \in \mathcal{B}$, puis comme $f$ est mesurable,
$f^{-1}(g^{-1}(C)) \in \mathcal{A}$.
\end{proof}

\section{Syst\`emes de Dynkin et lemme $\pi$-$\lambda$}

Le lemme $\pi$-$\lambda$ de Dynkin est un outil fondamental pour prouver que deux
mesures co\"incident.

\begin{definition}[$\pi$-syst\`eme et $\lambda$-syst\`eme]
\begin{itemize}
    \item Un \emph{$\pi$-syst\`eme} sur $X$ est une collection $\mathcal{P}
    \subset \mathcal{P}(X)$ stable par intersection finie.
    \item Un \emph{$\lambda$-syst\`eme} (ou \emph{syst\`eme de Dynkin}) sur $X$ est
    une collection $\mathcal{D} \subset \mathcal{P}(X)$ telle que~:
    \begin{enumerate}
        \item[(i)] $X \in \mathcal{D}$\,;
        \item[(ii)] si $A, B \in \mathcal{D}$ et $A \subset B$, alors
        $B \setminus A \in \mathcal{D}$\,;
        \item[(iii)] si $(A_n)$ est une suite croissante dans $\mathcal{D}$,
        alors $\bigcup_n A_n \in \mathcal{D}$.
    \end{enumerate}
\end{itemize}
\end{definition}

\begin{theoreme}[Lemme $\pi$-$\lambda$ de Dynkin]
\label{thm:pi_lambda}
Si $\mathcal{P}$ est un $\pi$-syst\`eme et $\mathcal{D}$ est un $\lambda$-syst\`eme
avec $\mathcal{P} \subset \mathcal{D}$, alors $\sigma(\mathcal{P}) \subset
\mathcal{D}$.
\end{theoreme}

\begin{proof}
Il suffit de montrer que $d(\mathcal{P})$, le plus petit $\lambda$-syst\`eme
contenant $\mathcal{P}$, est une $\sigma$-alg\`ebre (car alors
$\sigma(\mathcal{P}) \subset d(\mathcal{P}) \subset \mathcal{D}$).

Pour cela, il suffit de montrer que $d(\mathcal{P})$ est stable par intersection
finie (car un $\lambda$-syst\`eme stable par intersection finie est une
$\sigma$-alg\`ebre).

Pour $A \in d(\mathcal{P})$, posons
$\mathcal{G}_A = \{B \in d(\mathcal{P}) : A \cap B \in d(\mathcal{P})\}$.

\emph{\'Etape 1.} Si $A \in \mathcal{P}$, alors $\mathcal{G}_A$ est un
$\lambda$-syst\`eme contenant $\mathcal{P}$ (car $\mathcal{P}$ est un
$\pi$-syst\`eme), donc $\mathcal{G}_A \supset d(\mathcal{P})$.

\emph{\'Etape 2.} Pour tout $A \in d(\mathcal{P})$, $\mathcal{G}_A$ est un
$\lambda$-syst\`eme. Par l'\'etape 1, $\mathcal{P} \subset \mathcal{G}_A$ pour
tout $A \in d(\mathcal{P})$, donc $\mathcal{G}_A \supset d(\mathcal{P})$.

Ainsi $d(\mathcal{P})$ est stable par intersection finie, ce qui conclut.
\end{proof}

\begin{attention}[title=\textbf{Utilisation du lemme $\pi$-$\lambda$}]
Ce lemme est tr\`es utile pour montrer qu'une propri\'et\'e vraie sur un
$\pi$-syst\`eme s'\'etend \`a la $\sigma$-alg\`ebre engendr\'ee. Typiquement, pour
montrer que deux mesures $\mu$ et $\nu$ sur $\sigma(\mathcal{P})$ co\"incident, il
suffit de v\'erifier que $\mu = \nu$ sur le $\pi$-syst\`eme $\mathcal{P}$ (sous une
condition de $\sigma$-finitude).
\end{attention}

\section{Th\'eor\`eme d'unicit\'e des mesures}

\begin{theoreme}[Unicit\'e des mesures]
\label{thm:unicite_mesures}
Soient $\mu$ et $\nu$ deux mesures sur $(X, \sigma(\mathcal{P}))$, o\`u
$\mathcal{P}$ est un $\pi$-syst\`eme. Supposons qu'il existe une suite
$(E_n)_{n \geq 1} \subset \mathcal{P}$ avec $E_n \uparrow X$ et
$\mu(E_n) = \nu(E_n) < \infty$ pour tout $n$. Si $\mu = \nu$ sur $\mathcal{P}$,
alors $\mu = \nu$ sur $\sigma(\mathcal{P})$.
\end{theoreme}

\begin{proof}
Fixons $n$ et posons $\mu_n(\cdot) = \mu(\cdot \cap E_n)$ et
$\nu_n(\cdot) = \nu(\cdot \cap E_n)$. Ce sont des mesures finies.
La collection
\[
\mathcal{D}_n = \{A \in \sigma(\mathcal{P}) : \mu_n(A) = \nu_n(A)\}
\]
est un $\lambda$-syst\`eme (v\'erification directe). Par hypoth\`ese,
$\mathcal{P} \subset \mathcal{D}_n$. Par le lemme $\pi$-$\lambda$,
$\sigma(\mathcal{P}) \subset \mathcal{D}_n$, donc $\mu_n = \nu_n$ sur
$\sigma(\mathcal{P})$.

Pour tout $A \in \sigma(\mathcal{P})$, $A \cap E_n \uparrow A$, donc par
continuit\'e monotone~:
$\mu(A) = \lim_n \mu_n(A) = \lim_n \nu_n(A) = \nu(A)$.
\end{proof}

\section{Exercices}

\begin{exercice}
Soit $X$ un ensemble non d\'enombrable. Montrer que la collection
\[
\mathcal{A} = \{A \subset X : A \text{ est d\'enombrable ou } A^c \text{ est d\'enombrable}\}
\]
est une $\sigma$-alg\`ebre sur $X$.
\end{exercice}

\begin{exercice}
Montrer que $\mathcal{B}(\R)$ est engendr\'ee par les intervalles ouverts \`a
extr\'emit\'es rationnelles.
\end{exercice}

\begin{exercice}
Soit $f : \R \to \R$ une fonction continue. Montrer que $f$ est bor\'elienne.
\end{exercice}

\begin{exercice}
Montrer que $\mathcal{B}(\R)$ n'est \emph{pas} \'egale \`a $\mathcal{P}(\R)$.
\emph{Indication}~: utiliser un argument de cardinalit\'e ($\abs{\mathcal{B}(\R)} =
\abs{\R} = \mathfrak{c}$ alors que $\abs{\mathcal{P}(\R)} = 2^{\mathfrak{c}}$).
\end{exercice}

\begin{exercice}
Soient $(X, \mathcal{A})$ un espace mesurable et $f, g : X \to \R$ deux fonctions
mesurables. Montrer que $\{x \in X : f(x) < g(x)\} \in \mathcal{A}$.
\end{exercice}

\begin{exercice}[Atomes d'une $\sigma$-alg\`ebre]
Soit $(X, \mathcal{A})$ un espace mesurable. Pour $x \in X$, on d\'efinit
l'\emph{atome} de $x$ comme $[x] = \bigcap_{A \in \mathcal{A},\, x \in A} A$.
\begin{enumerate}
    \item Montrer que si $\mathcal{A}$ est engendr\'ee par une partition
    d\'enombrable, alors $[x] \in \mathcal{A}$ pour tout $x$.
    \item Donner un exemple o\`u $[x] \notin \mathcal{A}$.
\end{enumerate}
\end{exercice}
